\section{The canonical quotient of a forest marking}\label{the-canonical-quotient-of-a-forest-marking}

Suppose a group is generated by nonidentity involutions \(r_i\),
\(i\in I\). Its noncommutation graph has vertex set \(I\), and an edge
\(\{i,j\}\) precisely when \(r_i\) and \(r_j\) do not commute. We write
\(G_J=\langle r_j:j\in J\rangle\), as in the marking definitions. The intersection
property in this section always refers to the original group. It is not
assumed to survive an intermediate quotient.

\textbf{Lemma UA.Lift (descending lift).} Let \(G\) be a finite
nilpotent group, let \(A,B\le G\) satisfy \(A\cap B=1\), and fix an
integer \(m\ge1\). Put \(D_j=\gamma_j(G)\), \(A_j=A\cap D_j\),
\(B_j=B\cap D_j\). If the images of \(A_j,B_j\) generate \(D_j/D_{j+1}\)
for every \(j\ge m\), then \(D_j=A_jB_j\) for every \(j\ge m\).

\textbf{Proof.} Descend from a trivial term. Assuming
\(D_{j+1}=A_{j+1}B_{j+1}\), if \(a\in A_j,b\in B_j\) have equal images
in the next quotient, write \(ab^{-1}=a'b'\) with
\(a'\in A_{j+1},b'\in B_{j+1}\). Then \(a'^{-1}a=b'b\in A\cap B\), so
both images are trivial. Their images therefore form an internal direct
sum in the abelian quotient. Counting gives
\(|D_j|=|A_j||B_j|=|A_jB_j|\), proving equality as sets. No prior
normality or permutability of \(A,B\) is assumed. ∎

\textbf{Lemma UA.Support (graded support).} Let a positively graded Lie
algebra be generated in degree one by disjoint sets \(X,Y\), with
subalgebras \(L_X,L_Y\). If \([x,y]\in L_Y\) for \(x\in X,y\in Y\), then
\(L=L_X+L_Y\), degree by degree.

\textbf{Proof.} The derivation \(\operatorname{ad}x\) preserves \(L_Y\)
because it does so on its generators. Iterated brackets of those
derivations also preserve \(L_Y\), by Jacobi. Hence
\([L_X,L_Y]\subseteq L_Y\), and \(L_X+L_Y\) is a Lie subalgebra
containing all generators. The argument holds in characteristic two. ∎

\textbf{Theorem UA.Forest (forest quotient).} Suppose \(G\) is a finite
two-group generated by nonidentity involutions \(r_i\), with marked IP,
and its noncommuting pairs form a forest \(T\) on \(d\) vertices and
\(m\) edges. Then edge commutators are independent in
\(G'/\gamma_3(G)\), and \(G/\gamma_3(G)\) is the marked group \[
 (v,z)(w,t)=(v+w,z+t+\beta(v,w)),\quad
 v,w\in\mathbb F_2^d,\ z,t\in\mathbb F_2^m,               \tag{F0}
\] where, after orienting each edge, its component of \(\beta\) is
\(v_{\rm head}w_{\rm tail}\). Its order is \(2^{d+m}\).

\textbf{Proof.} IP gives irredundancy. UA.Phi gives \(G/G'=C_2^d\) and
\(G'\cap\langle r_i\rangle=1\). The lower-central graded Lie ring is
generated by these degree-one elements; all its graded components are
\(\mathbb F_2\)-spaces by bilinearity. Edge commutators span degree two.
If they satisfy a nonzero linear relation, root each forest component
and choose a participating edge whose child endpoint has maximal depth
in its component. Removing it separates the descendant set \(X\) from
all other vertices \(Y\), and no other participating edge is internal to
\(X\). The relation therefore writes the only nonzero cross bracket as a
sum of brackets internal to \(Y\). Lemma UA.Support shows every graded
component is generated by pure \(X\)- and pure \(Y\)-brackets. Their
representatives lie in \(G_X\cap\gamma_j(G)\) and
\(G_Y\cap\gamma_j(G)\).

Actual IP gives \(G_X\cap G_Y=1\), so Lemma UA.Lift gives
\(G'=(G'\cap G_X)(G'\cap G_Y)\). For the crossing endpoints
\(x\in X,y\in Y\), write \([r_x,r_y]=uv\) accordingly. Then
\(u\in G_X\cap G_{\{x\}\cup Y}=\langle r_x\rangle\), whence \(u=1\);
next \([r_x,r_y]\in G_Y\cap G_{\{x,y\}}=\langle r_y\rangle\), so it too
is one. This contradicts that the edge is noncommuting. Independence
follows.

Bilinearity makes (F0) associative; its coordinate markers are
involutions with independent central edge commutators. It is generated
by those markers, and its parabolic on any subset \(J\) consists exactly
of vertex coordinates supported on \(J\) and edge coordinates supported
on edges internal to \(J\). Thus it has IP. Restricting the vertex
coordinates to a subset and retaining exactly its internal edge
coordinates preserves the multiplication formula, so it also gives every
Boolean deletion. Collection in the corresponding class-two presentation
gives at most \(2^{d+m}\) elements, and this coordinate model attains
the bound. Independence in \(G'/\gamma_3(G)\) gives the same order for
\(G/\gamma_3(G)\); hence the marked map is an isomorphism. ∎
